\subsection{Special case: Binary Attack Detection}
\label{sec:binary}
We merge the five attack families into one Attack label and retain
Benign as the other class. A dedicated binary TM uses 400 clauses
per class (800 candidates), the same Boolean features, and the same
fit/selection/report partitions. Applying the selection-quarter pruning
rule with $\varepsilon=0.01$ removes $k^*=385$ clauses per class,
leaving 30 candidates, a 96.25\% reduction. The same CPSS grid selects
$(B,\pi_{\mathrm{thr}})=(25,0.9)$ for the full pool and $(25,0.8)$
for the pruned pool.

Table~\ref{tab:binary} follows the comparison in
Table~\ref{tab:headline}. Full-pool CAS raises binary macro-F1 from
$0.629\pm0.002$ to $0.823\pm0.016$ and accuracy from 0.842 to
0.885. Pruned CAS improves over the pruned TM in both metrics, but
falls below full-pool CAS. Preserving base-TM macro-F1 during pruning
therefore does not guarantee preservation of CAS performance.

\begin{table}[!htb]
\centering
\caption{Binary classification on the report quarter, mean$\pm$SD
over three seeds. CPSS settings are selected on the selection quarter.}
\label{tab:binary}
\footnotesize
\setlength{\tabcolsep}{3pt}
\begin{tabular}{@{}lccc@{}}
\toprule
Method & $(B,\pi_{\mathrm{thr}})$ & macro-F1 & Accuracy \\
\midrule
TM, full pool & --- & $0.629\pm0.002$ & $0.842\pm0.000$ \\
TM, pruned pool & --- & $0.690\pm0.034$ & $0.831\pm0.017$ \\
CAS, full pool & $(25,0.9)$ & $\mathbf{0.823\pm0.016}$ & $\mathbf{0.885\pm0.012}$ \\
CAS, pruned pool & $(25,0.8)$ & $0.756\pm0.007$ & $0.839\pm0.018$ \\
\bottomrule
\end{tabular}
\end{table}

Table~\ref{tab:binary-reduction} gives the binary counterpart of
Table~\ref{tab:reduction}, using matched CPSS settings for seed 42.
The Benign support shrinks from 178 to 19 clauses and the Attack
support from 178 to 18. Positive Attack evidence falls from 150 to
9 clauses. These counts describe compactness at a fixed setting;
Table~\ref{tab:binary} evaluates the separately selected operating points.

\begin{table}[!htb]
\centering
\caption{Binary signature sizes for seed 42 at matched settings
$B=15$, $\pi_{\mathrm{thr}}=0.8$.}
\label{tab:binary-reduction}
\footnotesize
\begin{tabular}{@{}lcccc@{}}
\toprule
& \multicolumn{2}{c}{Full (800)} & \multicolumn{2}{c}{Pruned (30)} \\
\cmidrule(lr){2-3}\cmidrule(lr){4-5}
Class & $|S|$ & $|S^+|$ & $|S|$ & $|S^+|$ \\
\midrule
Benign & 178 & 28 & 19 & 10 \\
Attack & 178 & 150 & 18 & 9 \\
\bottomrule
\end{tabular}
\end{table}

Full-pool CAS raises mean attack precision from 0.844 to 0.946,
while recall decreases from 0.987 to 0.910; attack F1 increases from
0.910 to 0.928. Its signed attack score achieves AUROC 0.938,
compared with 0.913 for the TM's raw attack score. Clipping the TM
score to $[0,T]$ and dividing by $T$ gives only 0.817, so the
unclipped score is the more informative ranking baseline. Pruned CAS
has AUROC 0.893, further illustrating the cost of aggressive compression.

\section{Discussion and Conclusion}
CAS turns a frozen TM clause pool into signed, inspectable class
evidence. For multiclass attribution, pruning retains 390 of 2,400
clauses and improves selected macro-F1 from 0.765 to 0.787.
Figure~\ref{fig:sweep}(a) shows that the pruned curves vary little
across stability thresholds, whereas the full-pool curves fluctuate
more. Figure~\ref{fig:sweep}(b) likewise shows little benefit from
increasing the number of complementary pairs for the pruned pool.
These fixed sensitivity slices support reduced parameter sensitivity
on this dataset; they do not establish statistical superiority at
every setting.

The binary comparison shows that compression has a task-dependent
cost: retaining only 30 candidates gives much shorter signatures but
lowers CAS macro-F1 by 0.068 relative to the full pool. Moreover,
full-pool binary CAS increases precision at the expense of recall,
raising the missed-attack fraction from 1.3\% to 9.0\%. Operating-point
selection must therefore reflect the intended false-alarm and
missed-attack costs, rather than macro-F1 alone.

The pruning margin condition concerns prediction preservation,
while the CPSS bound concerns low-selection-probability inclusions;
neither certifies causal explanations or analyst usefulness.
This retrospective evaluation covers three seeds and one capped
dataset, with prior test inspection and no matched XAI baseline.
Independent devices and time periods, semantic rule stability,
analyst evaluation, and measured runtime and energy are needed to
assess practical deployment. Overall, CAS improves attribution with
traceable signed evidence, while the full/pruned comparison exposes
the trade-off between signature size and predictive performance.

% All technical content and floats must finish by page 4.
% A subsequent page is reserved exclusively for references.
\clearpage

